\documentclass[11pt,a4paper]{article}
\usepackage[margin=23mm,headheight=14pt,headsep=8.5mm,footskip=10mm]{geometry}
\usepackage{amsmath,amssymb}
\usepackage{assets/tex/cm-bea}
\usepackage{xeCJK}
\setmainfont{Latin Modern Roman}
\setCJKmainfont{FandolSong-Regular.otf}[Path=assets/fonts/fandol/,BoldFont=FandolSong-Bold.otf,ItalicFont=FandolSong-Regular.otf]
\setCJKsansfont{FandolSong-Regular.otf}[Path=assets/fonts/fandol/,BoldFont=FandolSong-Bold.otf]
\xeCJKsetup{PunctStyle=quanjiao,CheckSingle=true}
\usepackage{setspace}\setstretch{1.12}
\setlength{\parindent}{2em}\setlength{\parskip}{3pt}
\usepackage{titlesec}
\titleformat{\section}{\centering\normalfont\large\bfseries}{}{0pt}{}
\titlespacing*{\section}{0pt}{1.5\baselineskip}{0.8\baselineskip}
\usepackage{needspace,enumitem,longtable,booktabs,array,calc}
\providecommand{\tightlist}{\setlength{\itemsep}{3pt}\setlength{\parskip}{0pt}}
\usepackage{fancyhdr}\pagestyle{fancy}\fancyhf{}
\fancyhead[L]{\small 数学手记}\fancyhead[R]{\small 楼红卫 B 组 · 第四十一批}
\fancyfoot[C]{\thepage}\renewcommand{\headrulewidth}{0.4pt}
\usepackage[hidelinks,unicode,linktoc=all]{hyperref}\usepackage{bookmark}
\hypersetup{pdftitle={楼红卫数学分析 B282—B288：Euler 积分、特殊函数与变分反例},pdfauthor={Vincent Arqain}}
\renewcommand{\contentsname}{题号目录}\setcounter{tocdepth}{1}\setcounter{secnumdepth}{0}
\widowpenalty=10000\clubpenalty=10000\displaywidowpenalty=10000
\predisplaypenalty=10000\interdisplaylinepenalty=10000
\raggedbottom\emergencystretch=2em
\begin{document}
\begin{center}
{\Large\bfseries 楼红卫《数学分析》B 组习题解答\par}
\vspace{7pt}{\large 第四十一批\quad B282—B288\par}
\vspace{5pt}{Euler 积分、特殊函数与变分反例\par}
\end{center}
本批收录习题 10.4.B 第 4—9 题及 10.5.B 第 1 题。正文依次给出原题、基础知识、思路来源、完整证明、易错条件与方法总结。数学公式采用 CM bea，中文采用 Fandol 宋体；公式编号与网页一致。
\begingroup\setlength{\parskip}{0pt}\tableofcontents\endgroup

\Needspace{9\baselineskip}
\section{问题 B282}\label{ux95eeux9898-b282}

\textbf{原题号：} 习题 10.4.B 第 4
题。\href{https://zj139.github.io/math-notes/files/problems/楼分析习题-v2.pdf\#page=32}{原题，PDF
第 32 页（印刷第 30 页）}

\textbf{题目}　计算

\begin{equation}
\lim_{x\to0^+}\sum_{n=1}^{\infty}\frac{nx^2}{1+n^3x^3}.
\tag{1}
\end{equation}

\Needspace{3\baselineskip}\textbf{基础知识}　紧区间上的 Riemann 和；可积尾部的一致估计；Euler 积分
\(B(a,b)=\Gamma(a)\Gamma(b)/\Gamma(a+b)\) 与反射公式。

\Needspace{3\baselineskip}\textbf{思路来源}　把分子写成 \(x(nx)\)，则网格宽度为 \(x\)，采样点为
\(nx\)。难点在积分区间无穷，紧区间上的 Riemann
和定理只能处理有限部分，还要独立控制尾和。

\Needspace{3\baselineskip}\textbf{解}　令 \(f(t)=t/(1+t^3)\)。对每个
\(x>0\)，题中级数收敛，因为其第 \(n\) 项不超过 \(1/(n^2x)\)。固定
\(A>0\)，记 \(N=\lfloor A/x\rfloor\)，则

\begin{equation}
x\sum_{n=1}^{N}f(nx)\longrightarrow\int_0^A f(t)\,dt
\quad(x\downarrow0).
\tag{2}
\end{equation}

这是 \([0,Nx]\) 上的右端点和，剩余区间 \([Nx,A]\) 的长度小于 \(x\)。当
\(0<x<A/2\) 时，\(N>A/(2x)\)，故

\begin{equation}
0\le x\sum_{n>N}f(nx)
\le\frac1x\sum_{n>N}\frac1{n^2}
\le\frac1{xN}\le\frac2A.
\tag{3}
\end{equation}

同时 \(\int_A^\infty f(t)\,dt\le1/A\)。由（2）、（3）先令
\(x\downarrow0\)，再令 \(A\to\infty\)，得到所求极限为
\(\int_0^\infty f\)。最后作换元 \(u=t^3\)：

\begin{equation}
\begin{aligned}
\int_0^\infty\frac{t}{1+t^3}\,dt
&=\frac13\int_0^\infty\frac{u^{2/3-1}}{1+u}\,du\\
&=\frac13 B\left(\frac23,\frac13\right)
=\frac{\pi}{3\sin(2\pi/3)}
=\frac{2\pi}{3\sqrt3}.
\end{aligned}
\tag{4}
\end{equation}

\Needspace{3\baselineskip}\textbf{易错条件}　不能把每个固定 \(n\)
的项的极限零与无限求和直接交换；主要贡献来自 \(n\) 与 \(1/x\)
同阶的范围。式（3）给出对小 \(x\) 一致的尾部控制，正是交换所需的依据。

\Needspace{3\baselineskip}\textbf{方法总结}　遇到 \(\sum x f(nx)\)，先识别网格，再把``有限区间
Riemann
和''和``无穷尾部一致小量''分开证明。积分可积并不自动保证任意采样和都有同样的尾估计。

\Needspace{9\baselineskip}
\section{问题 B283}\label{ux95eeux9898-b283}

\textbf{原题号：} 习题 10.4.B 第 5
题。\href{https://zj139.github.io/math-notes/files/problems/楼分析习题-v2.pdf\#page=32}{原题，PDF
第 32 页（印刷第 30 页）}

\textbf{题目}　计算

\begin{equation}
I=\int_0^\infty\frac{\log x}{x^2-1}\,dx.
\tag{5}
\end{equation}

\Needspace{3\baselineskip}\textbf{基础知识}　可去奇点；反常积分换元；非负级数的单调收敛；\(\sum_{n\ge1}n^{-2}=\pi^2/6\)。

\Needspace{3\baselineskip}\textbf{思路来源}　分母在 \(x=1\)
为零，但分子也同时为零，先检查这里是否真有发散。再用倒数换元把
\((1,\infty)\) 搬回
\((0,1)\)，此时几何级数的各项均为非负，可以安全换序。

\Needspace{3\baselineskip}\textbf{解}　当 \(x\to1\)
时，\(\log x/(x^2-1)\to1/2\)，故补上此值后连续。零端的绝对值为
\(O(|\log x|)\)，无穷端为
\(O(\log x/x^2)\)，所以（5）绝对收敛，无需主值解释。

在 \((1,\infty)\) 上令 \(x=1/t\)，得

\begin{equation}
\int_1^\infty\frac{\log x}{x^2-1}\,dx
=\int_0^1\frac{-\log t}{1-t^2}\,dt
=\int_0^1\frac{\log t}{t^2-1}\,dt.
\tag{6}
\end{equation}

于是利用非负项展开和单调收敛，

\begin{equation}
\begin{aligned}
I&=2\int_0^1\frac{-\log t}{1-t^2}\,dt
=2\sum_{k=0}^\infty\int_0^1t^{2k}(-\log t)\,dt\\
&=2\sum_{k=0}^\infty\frac1{(2k+1)^2}
=2\left(1-\frac14\right)\frac{\pi^2}{6}
=\frac{\pi^2}{4}.
\end{aligned}
\tag{7}
\end{equation}

其中
\(\int_0^1t^\alpha(-\log t)\,dt=(\alpha+1)^{-2}\)（\(\alpha>-1\)），可令
\(t=e^{-u}\) 后直接分部积分。

\Needspace{3\baselineskip}\textbf{易错条件}　\(x=1\) 是可去奇点。若把（5）拆成
\(\int\log x/(x-1)\) 与 \(\int\log x/(x+1)\)
两个全半轴积分，它们在无穷端分别发散，不能再作普通积分相减。

\Needspace{3\baselineskip}\textbf{方法总结}　倒数换元把两个端点的问题配对；几何级数则把对数积分转换为平方倒数和。先证明绝对收敛，再换元或合并，能避免无意义的发散相消。

\Needspace{9\baselineskip}
\section{问题 B284}\label{ux95eeux9898-b284}

\textbf{原题号：} 习题 10.4.B 第 6
题。\href{https://zj139.github.io/math-notes/files/problems/楼分析习题-v2.pdf\#page=32}{原题，PDF
第 32---33 页（印刷第 30---31 页）}

\textbf{题目}　设 \(p,q>1\) 为对偶数，即 \(1/p+1/q=1\)。原书假设
\(\sum|a_n|^p\) 和 \(\sum|b_n|^p\) 收敛，要求证明离散 Hardy--Hilbert
不等式

\begin{equation}
\sum_{m=1}^\infty\sum_{n=1}^\infty\frac{a_m b_n}{m+n}
\le\frac{\pi}{\sin(\pi/p)}
\left(\sum_{n=1}^\infty|a_n|^p\right)^{1/p}
\left(\sum_{n=1}^\infty|b_n|^q\right)^{1/q}.
\tag{8}
\end{equation}

\textbf{原题辨析}　完整教材第 249 页也把第二个假设印为 \(p\)，而右端使用
\(q\)。以下证明自然且有限的版本：\(a\in\ell^p\)、\(b\in\ell^q\)，并把左端加强为绝对值之和。若仅保留原印假设，当
\(1<p\le2\) 时由 \(\ell^p\subset\ell^q\) 仍可直接应用；当 \(p>2\)
时，原假设不能保证右端有限，甚至不能保证左端收敛，末尾给出反例。这里不把补充条件伪装成原题条件。

\Needspace{3\baselineskip}\textbf{基础知识}　Hölder
不等式；递减函数的积分比较；非负二重级数的单调收敛。Euler 积分给出

\begin{equation}
C_p:=\int_0^\infty\frac{t^{-1/q}}{1+t}\,dt
=B\left(\frac1p,\frac1q\right)
=\frac{\pi}{\sin(\pi/p)}
=\frac{\pi}{\sin(\pi/q)}.
\tag{9}
\end{equation}

\Needspace{3\baselineskip}\textbf{思路来源}　直接把核 \(1/(m+n)\) 平分给 Hölder
两因子，会留下发散的行和。应加入互相抵消的幂权
\((m/n)^{1/q}\)、\((n/m)^{1/p}\)，使每行的加权和恰由同一个 Euler
积分控制。

\Needspace{3\baselineskip}\textbf{证明}　固定 \(m\ge1\)，函数 \(t^{-1/q}/(m+t)\) 正且递减，因此

\begin{equation}
\begin{aligned}
\sum_{n=1}^\infty\frac{(m/n)^{1/q}}{m+n}
&\le m^{1/q}\int_0^\infty\frac{t^{-1/q}}{m+t}\,dt\\
&=\int_0^\infty\frac{u^{-1/q}}{1+u}\,du=C_p.
\end{aligned}
\tag{10}
\end{equation}

互换 \(m,n\) 及 \(p,q\)，同理有
\(\sum_{m\ge1}(n/m)^{1/p}/(m+n)\le C_p\)。对任意有限的 \(M,N\)，Hölder
给出

\begin{equation}
\begin{aligned}
\sum_{m=1}^M\sum_{n=1}^N\frac{|a_m b_n|}{m+n}
&=\sum_{m=1}^M\sum_{n=1}^N
\left(\frac{|a_m|^p(m/n)^{1/q}}{m+n}\right)^{1/p}
\left(\frac{|b_n|^q(n/m)^{1/p}}{m+n}\right)^{1/q}\\
&\le
\left(\sum_{m=1}^M|a_m|^p
\sum_{n=1}^N\frac{(m/n)^{1/q}}{m+n}\right)^{1/p}\\
&\quad\cdot
\left(\sum_{n=1}^N|b_n|^q
\sum_{m=1}^M\frac{(n/m)^{1/p}}{m+n}\right)^{1/q}\\
&\le C_p\|a\|_{\ell^p}\|b\|_{\ell^q}.
\end{aligned}
\tag{11}
\end{equation}

两个因子中的权经开方后分别为 \((m/n)^{1/(pq)}\) 与其倒数，故确实抵消。令
\(M,N\to\infty\)，单调收敛证明

\begin{equation}
\sum_{m,n\ge1}\frac{|a_m b_n|}{m+n}
\le C_p\|a\|_{\ell^p}\|b\|_{\ell^q}<\infty.
\tag{12}
\end{equation}

这同时证明二重级数绝对收敛、可任意换序，并推出实数列的（8）；复数列则对双重和取绝对值后同样成立。

为说明原印假设的不足，设 \(p>2\)，取 \(1/p<\alpha<1/2\)，令
\(a_n=b_n=n^{-\alpha}\)。两列都在 \(\ell^p\)，但对整数 \(N\ge1\)，

\begin{equation}
\sum_{m=N}^{2N}\sum_{n=N}^{2N}
\frac{m^{-\alpha}n^{-\alpha}}{m+n}
\ge\frac{N^2(2N)^{-2\alpha}}{4N}
=2^{-2\alpha-2}N^{1-2\alpha}\longrightarrow\infty.
\tag{13}
\end{equation}

故原印条件不能推出普通收敛的有限值结论。若允许非负扩充积分式约定，右端为无穷时仅得到平凡的扩充值比较。

\Needspace{3\baselineskip}\textbf{易错条件}　不能在尚未证明绝对收敛时直接重排有符号二重级数。也不能默默把假设中的
\(|b_n|^p\) 改成 \(|b_n|^q\) 而不说明；两指数只有在 \(p=q=2\)
时自动相同。

\Needspace{3\baselineskip}\textbf{方法总结}　齐次核搭配互相抵消的幂权，再用尺度换元把行和化为常数，是离散积分算子估计的常用方法。有限截断上的
Hölder 与非负单调极限，使收敛性和不等式一并得到。

\Needspace{9\baselineskip}
\section{问题 B285}\label{ux95eeux9898-b285}

\textbf{原题号：} 习题 10.4.B 第 7
题。\href{https://zj139.github.io/math-notes/files/problems/楼分析习题-v2.pdf\#page=33}{原题，PDF
第 33 页（印刷第 31 页）}

\textbf{题目}　计算 \(\psi\) 在 \(1/2,1/3,2/3,1/4,3/4,1/6,5/6\) 处的值。

\Needspace{3\baselineskip}\textbf{基础知识}　教材第 243 页定义
\(\psi(x)=(\log\Gamma(x))'=\Gamma'(x)/\Gamma(x)\)（\(x>0\)）。记 Euler
常数为 \(\gamma\)。使用 \(\Gamma\) 的反射公式、倍元公式与三倍元公式。

\Needspace{3\baselineskip}\textbf{思路来源}　同分母的一对特殊值，通常由``反射公式给差、倍元公式给和''联立求出。分母为六的值可由分母三的值再次二倍缩放得到，不必分别求积分。

\Needspace{3\baselineskip}\textbf{解}　由上一批 B280 已证明的
\(\int_0^\infty e^{-t}\log t\,dt=-\gamma\)，以及 \(\Gamma(1)=1\)，有

\begin{equation}
\psi(1)=-\gamma.
\tag{14}
\end{equation}

对反射公式和 \(m\) 倍元公式分别取对数导数，得到

\begin{equation}
\begin{aligned}
\psi(1-z)-\psi(z)&=\pi\cot(\pi z),&&0<z<1,\\
\sum_{k=0}^{m-1}\psi\left(z+\frac{k}{m}\right)
&=m\psi(mz)-m\log m,&&z>0,\quad m=2,3.
\end{aligned}
\tag{15}
\end{equation}

这里所用倍元公式为
\(\prod_{k=0}^{m-1}\Gamma(z+k/m)=(2\pi)^{(m-1)/2}m^{1/2-mz}\Gamma(mz)\)，故（15）中的
\(-m\log m\) 不能省略。

在（15）的第二式中先取 \((m,z)=(2,1/2)\)，再取 \((3,1/3)\)，得到

\begin{equation}
\begin{aligned}
\psi(1/2)&=-\gamma-2\log2,\\
\psi(1/3)+\psi(2/3)&=-2\gamma-3\log3.
\end{aligned}
\tag{16}
\end{equation}

又由反射关系
\(\psi(2/3)-\psi(1/3)=\pi/\sqrt3\)，可得分母三的一对值。类似地，

\begin{equation}
\begin{aligned}
\psi(1/4)+\psi(3/4)&=2\psi(1/2)-2\log2,\\
\psi(3/4)-\psi(1/4)&=\pi,\\
\psi(1/6)+\psi(2/3)&=2\psi(1/3)-2\log2,\\
\psi(5/6)-\psi(1/6)&=\sqrt3\,\pi.
\end{aligned}
\tag{17}
\end{equation}

解这些二元一次关系，七个结果为

\begin{equation}
\begin{aligned}
\psi(1/2)&=-\gamma-2\log2,\\
\psi(1/3)&=-\gamma-\frac32\log3-\frac{\pi}{2\sqrt3},\\
\psi(2/3)&=-\gamma-\frac32\log3+\frac{\pi}{2\sqrt3},\\
\psi(1/4)&=-\gamma-3\log2-\frac\pi2,\\
\psi(3/4)&=-\gamma-3\log2+\frac\pi2,\\
\psi(1/6)&=-\gamma-2\log2-\frac32\log3-\frac{\sqrt3\,\pi}{2},\\
\psi(5/6)&=-\gamma-2\log2-\frac32\log3+\frac{\sqrt3\,\pi}{2}.
\end{aligned}
\tag{18}
\end{equation}

\Needspace{3\baselineskip}\textbf{易错条件}　反射式左端是
\(\psi(1-z)-\psi(z)\)，顺序反过来就要改符号。\(\psi\) 是 \(\log\Gamma\)
的导数，不能把 \(\Gamma\) 的特殊值直接取对数当作答案。

\Needspace{3\baselineskip}\textbf{方法总结}　先确定一个基准值，再用函数方程组织未知量。相同分母用和差关系，分母倍增用倍元关系，可以把特殊函数值的计算化为少量线性代数运算。

\Needspace{9\baselineskip}
\section{问题 B286}\label{ux95eeux9898-b286}

\textbf{原题号：} 习题 10.4.B 第 8
题。\href{https://zj139.github.io/math-notes/files/problems/楼分析习题-v2.pdf\#page=33}{原题，PDF
第 33 页（印刷第 31 页）}

\textbf{题目}　证明 \(\Gamma\) 函数在定义域内解析。

\Needspace{3\baselineskip}\textbf{基础知识}　实解析函数的局部幂级数定义；绝对可积条件下的求和积分交换；\(\Gamma(x+1)=x\Gamma(x)\)。教材第
236 页先在 \(x>0\) 上以 Euler
积分定义，再用递推扩充到除零与负整数外的所有实数。

\Needspace{3\baselineskip}\textbf{思路来源}　任意阶可微仍不足以推出解析。应直接把
\(t^{x+h-1}=t^{x-1}e^{h\log t}\)
展开，并证明所有绝对值之和有可积控制；得到的才是真正收敛到函数本身的
Taylor 级数。

\Needspace{3\baselineskip}\textbf{证明}　固定 \(x_0>0\)，任取 \(0<r<x_0\)。对 \(|h|\le r\)，有

\begin{equation}
\begin{aligned}
&t^{x_0-1}e^{-t}\sum_{k=0}^\infty
\frac{|h\log t|^k}{k!}
=t^{x_0-1}e^{-t}e^{|h||\log t|}\\
&\qquad\le
\begin{cases}
t^{x_0-r-1},&0<t<1,\\
t^{x_0+r-1}e^{-t},&t\ge1.
\end{cases}
\end{aligned}
\tag{19}
\end{equation}

右端在 \((0,\infty)\) 可积。于是可以逐项积分，得到

\begin{equation}
\begin{aligned}
\Gamma(x_0+h)
&=\int_0^\infty t^{x_0-1}e^{-t}e^{h\log t}\,dt\\
&=\sum_{k=0}^\infty c_k h^k,\qquad |h|<x_0,\\
c_k&=\frac1{k!}\int_0^\infty t^{x_0-1}e^{-t}(\log t)^k\,dt.
\end{aligned}
\tag{20}
\end{equation}

式（19）还说明
\(\sum|c_k|r^k<\infty\)，故（20）是局部一致收敛的幂级数，并且
\(c_k=\Gamma^{(k)}(x_0)/k!\)。这证明正半轴上的实解析性。

若 \(x_0<0\) 且 \(x_0\) 不是整数，选整数 \(N\) 使 \(x_0+N>0\)。在
\(x_0\) 的一个足够小的邻域内，

\begin{equation}
\Gamma(x)=\frac{\Gamma(x+N)}
{x(x+1)\cdots(x+N-1)}.
\tag{21}
\end{equation}

分子解析，分母为非零多项式，所以商也解析。因此 \(\Gamma\)
在教材的全部实定义域 \(\mathbb R\setminus\{0,-1,-2,\ldots\}\)
上解析。事实上，在（19）中允许复数
\(h\)，同一绝对值估计也直接证明右半平面的全纯性；递推则给出避开这些极点的全纯延拓。

\Needspace{3\baselineskip}\textbf{易错条件}　只证明所有 \(\Gamma^{(k)}\)
存在，不能代替（20）。负实轴上不能沿用零端发散的 Euler
积分，也不能到处写 \(\log\Gamma(x)\)；这里使用非零分母的商来延拓。

\Needspace{3\baselineskip}\textbf{方法总结}　含参积分的解析性可由``指数展开＋绝对值级数的可积控制''直接建立。先在积分收敛域证明，再用函数方程扩大定义域，能清楚区分积分表达式与函数本身的定义范围。

\Needspace{9\baselineskip}
\section{问题 B287}\label{ux95eeux9898-b287}

\textbf{原题号：} 习题 10.4.B 第 9
题。\href{https://zj139.github.io/math-notes/files/problems/楼分析习题-v2.pdf\#page=33}{原题，PDF
第 33 页（印刷第 31 页）}

\textbf{题目}　证明例 9.4.7 中的结论。

\textbf{引用内容核实}　完整教材第 181---182 页（扫描 PDF 第 195---196
页）的例 9.4.7 为：设

\begin{equation}
a_0=0,\qquad a_1=\frac23,\qquad
a_{n+1}=\frac{2}{3n}a_n+a_{n-1}\quad(n\ge1).
\tag{22}
\end{equation}

教材先求得母函数，随后留下积分表示（原式 9.4.18）及渐近等价式（原式
9.4.19）的证明：

\begin{equation}
\begin{aligned}
a_n&=\frac{n}{\Gamma(2/3)\Gamma(1/3)}
\int_{-1}^1\left(\frac{1+s}{1-s}\right)^{1/3}s^{n-1}\,ds,\quad n\ge1,\\
a_n&\sim\frac{2^{1/3}}{\Gamma(1/3)}\,n^{1/3},
\qquad n\to\infty.
\end{aligned}
\tag{23}
\end{equation}

以下一并补证母函数、积分表示与精确渐近常数；不是只给出 \(a_n\)
的数量级。

\Needspace{3\baselineskip}\textbf{基础知识}　幂级数母函数；端点可积权的分部积分；Beta
积分；控制收敛。记

\begin{equation}
C=\Gamma(2/3)\Gamma(1/3),\qquad
w(s)=\left(\frac{1+s}{1-s}\right)^{1/3}\quad(-1<s<1).
\tag{24}
\end{equation}

\Needspace{3\baselineskip}\textbf{思路来源}　分母为 \(n\)
的递推经母函数求导可化为微分方程。要证明积分表示，则反过来验证矩积分满足同一个递推；渐近行为由权在
\(s=1\) 的可积奇性决定，同时必须排除 \(s=-1\) 的交错贡献。

\Needspace{3\baselineskip}\textbf{证明}　由（22）归纳可得 \(0\le a_n\le n\)。其中归纳步为
\(a_{n+1}\le2/3+n-1\le n+1\)，初值显然满足。因此
\(S(x)=\sum_{n\ge1}a_nx^{n-1}\) 在 \(|x|<1\) 收敛，递推给出

\begin{equation}
\begin{aligned}
\sum_{n\ge1}\frac{2a_n}{3n}x^n
&=(1-x^2)S(x)-a_1,\\
\big((1-x^2)S(x)\big)'&=\frac23S(x),\\
S(x)&=\frac23(1+x)^{-2/3}(1-x)^{-4/3}.
\end{aligned}
\tag{25}
\end{equation}

最后一式由 \(S(0)=2/3\) 解一阶微分方程得到，正是教材原式（9.4.17）。

为证（23）的积分公式，令
\(M_k=\int_{-1}^1w(s)s^k\,ds\)（\(k\ge0\)），并令
\(c_0=0\)、\(c_n=nM_{n-1}/C\)。由于
\((1-s^2)w(s)=(1+s)^{4/3}(1-s)^{2/3}\)
在两端趋零，且导数可积，可在截断区间分部积分后令端点趋于 \(\pm1\)，得到

\begin{equation}
\begin{aligned}
0&=\int_{-1}^1
\frac{d}{ds}\big(s^k(1-s^2)w(s)\big)\,ds\\
&=kM_{k-1}+\frac23M_k-(k+2)M_{k+1}\quad(k\ge1),\\
2M_1&=\frac23M_0\quad(k=0).
\end{aligned}
\tag{26}
\end{equation}

这里用了 \(w'/w=2/[3(1-s^2)]\)。令 \(u=(1+s)/2\)，Beta 积分给出

\begin{equation}
M_0=2B\left(\frac43,\frac23\right)
=\frac23C,\qquad c_1=\frac23.
\tag{27}
\end{equation}

由（26）得 \(c_2=(2/3)c_1\)，且对 \(n\ge2\)，

\begin{equation}
(n+1)M_n=\frac23M_{n-1}+(n-1)M_{n-2},
\qquad c_{n+1}=\frac{2}{3n}c_n+c_{n-1}.
\tag{28}
\end{equation}

所以 \(c_n\) 与 \(a_n\) 初值、递推完全相同，必有
\(c_n=a_n\)，积分表示得证。

为求渐近式，把 \(M_{n-1}\) 分成正、负半区间。正半区间令 \(t=n(1-s)\)，有

\begin{equation}
\begin{aligned}
P_n&:=\int_0^1(1+s)^{1/3}(1-s)^{-1/3}s^{n-1}\,ds,\\
n^{2/3}P_n
&=\int_0^n(2-t/n)^{1/3}t^{-1/3}(1-t/n)^{n-1}\,dt.
\end{aligned}
\tag{29}
\end{equation}

把右端被积函数在 \(t\ge n\) 处补零。对 \(n\ge2\)，
\((1-t/n)^{n-1}\le e^{-t/2}\)，故它受可积函数
\(2^{1/3}t^{-1/3}e^{-t/2}\) 控制。逐点极限为
\(2^{1/3}t^{-1/3}e^{-t}\)，于是

\begin{equation}
n^{2/3}P_n\longrightarrow2^{1/3}\Gamma(2/3).
\tag{30}
\end{equation}

负半区间令 \(s=-u\)，则当 \(n\ge2\) 时，

\begin{equation}
\begin{aligned}
\left|\int_{-1}^0w(s)s^{n-1}\,ds\right|
&\le\int_0^1(1-u)^{1/3}u^{n-1}\,du\\
&\le\int_0^1v^{1/3}e^{-(n-1)v}\,dv
\le\frac{\Gamma(4/3)}{(n-1)^{4/3}}.
\end{aligned}
\tag{31}
\end{equation}

它是 \(o(n^{-2/3})\)，不会改变正半区间的主项。代入积分表示，得到

\begin{equation}
a_n=\frac nC\big(P_n+o(n^{-2/3})\big)
\sim\frac{2^{1/3}}{\Gamma(1/3)}\,n^{1/3}.
\tag{32}
\end{equation}

\Needspace{3\baselineskip}\textbf{易错条件}　原例的递推系数是 \(2/(3n)\)，积分前系数是
\(n/C\)。\(s^{n-1}\)
在负半轴的符号随奇偶改变，不能未经估计就把整个积分的主项只归于
\(s=1\)。式（31）补足了这一点。

\Needspace{3\baselineskip}\textbf{方法总结}　母函数负责识别结构，矩积分通过递推唯一性验证精确公式，端点缩放负责提取渐近常数。这三种工具各自解决不同问题，避免把母函数的奇点形状当作未经证明的系数渐近定理。

\Needspace{9\baselineskip}
\section{问题 B288}\label{ux95eeux9898-b288}

\textbf{原题号：} 习题 10.5.B 第 1
题。\href{https://zj139.github.io/math-notes/files/problems/楼分析习题-v2.pdf\#page=33}{原题，PDF
第 33 页（印刷第 31 页）}

\textbf{题目}　举例说明，对于
\(A\in C([a,b];\mathbb R^{(2m)\times(2m)})\)，\(A(x)^T=A(x)\)，条件

\begin{equation}
\int_a^b
\begin{pmatrix}\varphi(x)\\\varphi'(x)\end{pmatrix}^{T}
A(x)
\begin{pmatrix}\varphi(x)\\\varphi'(x)\end{pmatrix}\,dx
\ge0,\qquad
\forall\varphi\in C_0^1([a,b];\mathbb R^m)
\tag{33}
\end{equation}

并不蕴涵 \(A(x)\ge0\) 对所有 \(x\in[a,b]\) 成立。这里
\(a<b\)，\(m\ge1\)，\(A\ge0\) 指半正定；检验函数满足两端为零。

\Needspace{3\baselineskip}\textbf{基础知识}　实对称矩阵的二次型；分块矩阵；全导数的积分；零边界条件。

\Needspace{3\baselineskip}\textbf{思路来源}　检验向量被限制为
\((\varphi,\varphi')\)，其上下两部分不能独立选取。如果能让被积式恰为一个全导数，零边界就会令所有积分消失，而矩阵仍可以是不定的。

\Needspace{3\baselineskip}\textbf{证明}　取常矩阵

\begin{equation}
A(x)\equiv
\begin{pmatrix}0&I_m\\I_m&0\end{pmatrix}.
\tag{34}
\end{equation}

它连续、实对称。对任意允许的检验函数，

\begin{equation}
\begin{aligned}
\int_a^b
\begin{pmatrix}\varphi\\\varphi'\end{pmatrix}^{T}
A
\begin{pmatrix}\varphi\\\varphi'\end{pmatrix}\,dx
&=2\int_a^b\varphi(x)^T\varphi'(x)\,dx\\
&=\big[|\varphi(x)|^2\big]_a^b=0.
\end{aligned}
\tag{35}
\end{equation}

所以（33）成立。但任取非零的 \(u\in\mathbb R^m\)，有

\begin{equation}
\begin{pmatrix}u\\-u\end{pmatrix}^{T}
A
\begin{pmatrix}u\\-u\end{pmatrix}
=-2|u|^2<0.
\tag{36}
\end{equation}

故 \(A(x)\) 在每一点都不是半正定矩阵。实际上它的特征值为
\(1,-1\)，各重数 \(m\)；此反例对任意区间长度及任意 \(m\ge1\) 都适用。

\Needspace{3\baselineskip}\textbf{易错条件}　不能由对所有 \((\varphi,\varphi')\)
的积分检验，推出对每点所有独立向量 \((u,v)\) 的二次型都非负。试图把
\(\varphi\) 局部化为常数时，截断区域仍会产生导数项，不能丢掉。

\Needspace{3\baselineskip}\textbf{方法总结}　全导数项在零边界的变分积分中不可见，却能改变点态二次型。这一例子说明：积分意义的非负性受微分约束与边界条件共同影响，不能直接套用对任意独立检验向量成立的点态判别。

\Needspace{9\baselineskip}
\section{七题方法总结}\label{ux4e03ux9898ux65b9ux6cd5ux603bux7ed3}

\begin{longtable}[]{@{}
  >{\raggedright\arraybackslash}p{(\columnwidth - 4\tabcolsep) * \real{0.12}}
  >{\raggedright\arraybackslash}p{(\columnwidth - 4\tabcolsep) * \real{0.44}}
  >{\raggedright\arraybackslash}p{(\columnwidth - 4\tabcolsep) * \real{0.44}}@{}}
\toprule\noalign{}
\begin{minipage}[b]{\linewidth}\raggedright
题号
\end{minipage} & \begin{minipage}[b]{\linewidth}\raggedright
核心结果
\end{minipage} & \begin{minipage}[b]{\linewidth}\raggedright
关键方法
\end{minipage} \\
\midrule\noalign{}
\endhead
\bottomrule\noalign{}
\endlastfoot
B282 & 极限为 \(2\pi/(3\sqrt3)\) & 紧区间 Riemann 和、统一尾界、Euler
积分 \\
B283 & 对数积分为 \(\pi^2/4\) & 可去奇点、倒数换元、非负级数换序 \\
B284 & 离散 Hardy--Hilbert 的对偶指数版本 & 抵消幂权、积分比较、有限截断
Hölder \\
B285 & 七个有理点的 \(\psi\) 值 & 基准值、反射给差、倍元给和 \\
B286 & \(\Gamma\) 在全部实定义域上解析 & 指数展开的可积控制、递推延拓 \\
B287 & 原例积分表示及精确的 \(n^{1/3}\) 主项 &
母函数、矩递推、两端分别缩放估计 \\
B288 & 非负积分二次型的常矩阵反例 & 全导数、零边界、受约束检验向量 \\
\end{longtable}

\end{document}
